\section{Semantics}

%Each variable $x$ is associated with a set of expressions $E_x$.
%\subsection{Exogenous Variables}
%\subsection{Endogenous Variables}
Let $\atoms$ and $\variables$ be disjoint sets of ground terms called \emph{atoms} and \emph{variables}. %the Herbrand universe.
Let $\discourse$ be a set of \emph{values} called \emph{universe of discourse}.
We will assume that the universe of discourse contains two distinguished \emph{boolean} values, $\True$ and $\False$. % values are in the universe of discourse.
A \emph{valuation} $\vec{v}$ or ``assignment'' or ``interpretation'' is a (partial) mapping from terms in $\atoms \cup \variables$ to values in $\discourse$, i.e., $\vec{v}: \atoms \cup\variables \partialto \discourse$.
For any $x \in \atoms \cup \variables$, writing $\vec{v}(x)$ implicitly implies that $\vec{v}$ is defined on $x$.
We call \emph{theory part} of $\vec{v}$ the mapping $\theoryPart{\vec{v}}: \variables \partialto \discourse$.
With a small abuse of notation, we call \emph{model part} of $\vec{v}$ the set of atoms $\modelPart{\vec{v}} = \{ x \suchthat x \in \atoms, \vec{v}(x) = \True \}$.

A valuation $\vec{u}$ is smaller than a valuation $\vec{v}$ if $\vec{u}$ and $\vec{v}$ have the same model part and they coincide on every variable in the domain of $\vec{u}$.
That is, for every term $y \in \variables$ and value $a \in \discourse$, $y \mapsto a \in \vec{u}$ implies that $y \mapsto a \in \vec{v}$.
Furthermore, $\vec{u}$ is strictly smaller than $\vec{v}$ is $\vec{u}$ is smaller than $\vec{v}$ and there exists a term $y$ and a vaue $a$ such that $y \mapsto a \in \vec{v}$ while $y \mapsto a \notin \vec{u}$.

%We call terms in the domain of a valuation \emph{variables}.

\subsection{Evaluation of expressions}

%\begin{align*}
%\evaluae{\vec{v}}{x} & =\\
%\end{align*}

\inference[constants]{}{\vec{v}\entails\redto{\texttt{val(t,c)}}{c}}

\bigskip
\inference[variable]{}{\vec{v}\entails\redto{\texttt{variable($x$)}}{\vec{v}(x)}}

\bigskip
\inference[tuple]{\vec{v}\entails\redto{e_1}{a_1}&\vec{v}\entails\redto{e_2}{a_2}&\dots&\vec{v}\entails\redto{e_k}{a_k}&}{\vec{v}\entails\redto{(e_1,\dots,e_k)}{(a_1,\dots,a_k)}}

\bigskip
\inference[application]{\vec{v}\entails\redto{e_o}{f}&\vec{v}\entails\redto{e_1}{a_1}&\vec{v}\entails\redto{e_2}{a_2}&\dots&\vec{v}\entails\redto{e_k}{a_k}&}{\vec{v}\entails\redto{\texttt{operation($e_o$,[$e_1$,$\dots$,$e_k$])}}{f(a_1,\dots,a_k)}}

\subsection{Evaluation of statements}

\inference[assert]{\vec{v}\entails\redto{e}{\True}}{\vec{v}\entails\redto{\texttt{assert($e$)}}{\vec{v}}}
\bigskip

\inference[assign]{\vec{v}\entails\redto{e}{a}}{\vec{v}\entails\redto{\texttt{assign($x$,$e$)}}{\vec{w}}}
where $\vec{w} = \{ x \mapsto a \} \cup \{ z \mapsto \vec{v}(z) \suchthat z \in \variables \text{ and } z \neq x\}$

\bigskip
\inference[if then branch]{\vec{v}\entails\redto{e}{\True}&\vec{v}\entails\redto{s_1}{\vec{w}}}{\vec{v}\entails\redto{\texttt{if($e$,$s_1$,$s_2$)}}{\vec{w}}}

\bigskip
\inference[if else branch]{\vec{v}\entails\redto{e}{\False}&\vec{v}\entails\redto{s_2}{\vec{w}}}{\vec{v}\entails\redto{\texttt{if($e$,$s_1$,$s_2$)}}{\vec{w}}}

\bigskip
\inference[noop]{}{\vec{v}\entails\redto{\texttt{noop}}{\vec{v}}}

\bigskip
\inference[seq2]{\vec{v}\entails\redto{s_1}{\vec{u}}&\vec{u}\entails\redto{s_2}{\vec{w}}}{\vec{v}\entails\redto{\texttt{seq2($s_1$,$s_2$)}}{\vec{w}}}

\bigskip
\inference[while body]{\vec{v}\entails\redto{e}{\True}&\vec{v}\entails\redto{s}{\vec{u}}&\vec{u}\entails\redto{\texttt{while($e$,$s$)}}{\vec{w}}}{\vec{v}\entails\redto{\texttt{while($e$,$s$)}}{\vec{w}}}

\bigskip
\inference[while exit]{\vec{v}\entails\redto{e}{\False}}{\vec{v}\entails\redto{\texttt{while($e$,$s$)}}{\vec{v}}}

\bigskip
\inference[bounded while body]{\vec{v}\entails\redto{e}{\True}&\vec{v}\entails\redto{s}{\vec{u}}&\vec{u}\entails\redto{\texttt{while($n-1$,$e$,$s$)}}{\vec{w}}}{\vec{v}\entails\redto{\texttt{while($n$,$e$,$s$)}}{\vec{w}}}

\bigskip
\inference[bounded while exit 1]{}{\vec{v}\entails\redto{\texttt{while($0$,$e$,$s$)}}{\vec{v}}}

\bigskip
\inference[bounded while exit 2]{\vec{v}\entails\redto{e}{\False}}{\vec{v}\entails\redto{\texttt{while($n$,$e$,$s$)}}{\vec{v}}}

TODO python

\subsection{Extended model}
Let $\vec{v}$ be a valuation and $\modelPart{\vec{v}}$ be its model part.
We say that $\vec{v}$ is \emph{consistent} if the following conditions hold.
\begin{itemize}
\item for every atom $\texttt{variable_define($y$,$e$)} \in \modelPart{\vec{v}}$, $e$ reduces to some value and $y$ maps to it, that is $\vec{v}\entails\redto{e}{\vec{v}(a)}$;
\item for every atom $\texttt{variable_declare($y$,boolDomain)} \in \modelPart{\vec{v}}$, $y$ is mapped to a boolean value, that is $\vec{v}(y) \in \mathbb{B}$;
\item for every atom $\texttt{variable_declare($y$,fromFacts)} \in \modelPart{\vec{v}}$, there exists an expression $e$ such that $\texttt{variable_domain($y$,$e$)} \in \modelPart{\vec{v}}$ and $\vec{v}\entails\redto{e}{\vec{v}(y)}$;
\item for every atom $\texttt{execution_declare($p$,$s$,$i$,$o$)} \in \modelPart{\vec{v}}$ such that $\texttt{execution_run($p$)} \in \modelPart{\vec{v}}$, the valuations $\vec{i} = \{ x \mapsto \vec{v}(\texttt{execution_input($s$,$x$)}) \suchthat x \in i \}$ and $\vec{o} = \{ x \mapsto \vec{v}(\texttt{execution_output($s$,$x$)}) \suchthat x \in o \}$ satisfy $\vec{i}\entails\redto{s}{\vec{o}}$; and
\item for every atom $\texttt{ensure($e$)} \in \modelPart{\vec{v}}$, we have $\vec{v}\entails\redto{e}{\True}$.
\end{itemize}

Let $P$ be an answer set program, $\vec{v}$ be a valuation, and $\modelPart{\vec{v}}$ its model part.
We say that $\vec{v}$ is a \emph{solution} of $P$, and write $\vec{v} \models P$, whenever the following conditions hold.

\begin{itemize}
\item $\modelPart{\vec{v}}$ is an answer set of $P$;
\item $\vec{v}$ is consistent; and
\item There does not exist any consistent valuation $\vec{u}$ strictly smaller than $\vec{v}$.
\end{itemize}


\section{Typing}
\section{Reasoning modes}
\subsection{Extended brave and cautious models}
Let $P$ be an answer set program that admits at least one solution.
The \emph{extended brave model} of $P$ is the relation $\brave{P}$ over $(\atoms \cup \variables) \times \discourse$ defined so that a term $x$ is in relation to a value $a$ whenever $x$ is mapped to $a$ in some solution of $P$.
Similarly, the \emph{extended cautious model} of $P$ is the relation $\cautious{P}$ defined so that $x$ is in relation to $a$ whenever $x$ is mapped to $a$ in all solutions of $P$.
In other words, the extended brave model of $P$ is the relation $\brave{P} = \bigcup_{\vec{v}\models P} \vec{v}$ and the extended cautious model of $P$ is the relation $\cautious{P} = \bigcap_{\vec{v}\models P} \vec{v}$.

\subsection{Supremum and infimum models}
%Let $P$ be an answer set program that admits at least one solution.
The \emph{supremum model} of $P$ is the mapping $\supremum{P} \colon (\atoms \cup \variables) \rightarrow \discourse$ where every term is mapped to the least upper bound of values $x$ takes across solutions of $P$.
Similarly, the \emph{infimum model} of $P$ is the relation $\infimum{P} \colon (\atoms \cup \variables) \rightarrow \discourse$ where every term is mapped to the greatest lower bound of values $x$ takes across solutions of $P$. %defined so that $x$ is in relation to $a$ whenever $x$ is mapped to $a$ in all solutions of $P$.
In other words, the supremum model of $P$ is the mapping defined by $\supremum{P}(x) \mapsto \bigvee_{\vec{v}\models P} \vec{v}(x)$ and the infimum model of $P$ is defined by $\infimum{P}(x) \mapsto \bigwedge_{\vec{v}\models P} \vec{v}(x)$.

\begin{fact}
If $\discourse$ is a complete lattice, then the supremum and infimum models are guaranteed to exist.
\end{fact}
\begin{fact}
If ther are no variables, $\variables = \emptyset$, then the extended brave model and the supremum coincide with the brave model $P$, and the extended cautious and infimum model coincide with the cautious model of $P$.
\end{fact}

\section{Implementation}
The dependencies of $x$

dependency graph

Acyclic structural equations
